%\clearpage

\section {Introduction}
\label{sec:intro}

In this technical note, we report the results of a 3-flavour $\bar{\nu}_{e}$-appearance analysis,
performed by the {\em VALOR} group,        
on the combined Run 5+6 dataset, using only data taken in antineutrino beam mode (
\TotalPot{} POT
).
This dataset was collected in two run periods:
May - June 2014 (Run 5, \RunVPot{} POT in antineutrino beam mode) and November 2014 - June 2015 
(Run 6, 
\RunVIPot{} 
POT in antineutrino beam mode).
This analysis builds on our other analyses presented in Refs. 
\cite{OurRun1DisappAnalysis},
\cite{OurRun12DisappAnalysis}, 
\cite{OurRun12DisappAnalysis3Flavour},
\cite{OurRun123DisappAnalysis3Flavour},
\cite{OurRun1234aDisappAnalysis3Flavour},
\cite{OurRun1234DisappAnalysis3Flavour},
\cite{OurRun123JointOscillationAnalysis},
\cite{OurRun1234JointOscillationAnalysis} and
\cite{OurRun56AntiNumuDisappearanceAnalysis3Flavour}
, and adapts these analyses  
to antineutrino beam mode data.
The main difference is that in the past a fully frequentist analysis was performed.
Now a hybrid frequentist-bayesian method \cite{CousinsHighland}, 
which consists of marginalizing instead of fitting 
the systematic parameters, is used.
Prior distributions obtained from the T2K+reactor joint fit 
\cite{Mach3JointFitRun14} were used in the analysis.

The goal of this analysis is to search for $\bar{\nu}_e$ appearance in the PMNS matrix contest
and not to search for any physics beyond the Standard Model.
For this reason, differently from the $\bar{\nu}_{\mu}$-disappearance analysis \cite{OurRun56AntiNumuDisappearanceAnalysis3Flavour}
common oscillation parameters are used for $\nu$ and $\bar{\nu}$.
The consistency of the data with the expectation is given by comparing the predicted reconstructed energy spectrum of 1 $e$-like ring 
events at SuperK to the observed spectrum. The method of calculation of this predicted spectrum is presented in Sec.~\ref{sec:spectrum_prediction}.
This analysis uses flux tuning based on recent NA61 data \cite{FluxTuning}\cite{NA61T2K}. 
It also uses an indirect extrapolation from ND280 to SuperK, which is implemented via a fit of flux and cross-section parameters to
ND280 measurements of the muon momentum-angle distributions of $\bar{\nu}_\mu$ and  $\nu_\mu$
CC0$\pi$-like, CC1$\pi$-like and CC other events whose 
vertices are in FGD1\cite{BANFF2015}.
Updated nominal predictions of the reconstructed energy spectrum of single $e$-like ring events in SuperK are presented for Run 5+6
in Sec.~\ref{sec:spectrum_prediction}

The oscillation analysis is performed using neutrino oscillation probabilities computed in a 3-flavour
framework including constant-density matter effects. 
Details of our oscillation probability calculations, estimates of the numerical accuracy of our
calculations and comparisons with calculations used by the alternative T2K oscillation analyses 
are given in Appendices A and B of Ref. \cite{OurRun123DisappAnalysis3Flavour}.

As in our previous analyses, the oscillation fit is made using the likelihood-ratio method \cite{StatAnalysisCowan}.
The mathematical details of this method are given in Sec.\ 3.2 in Ref.\ \cite{OurRun1DisappAnalysis}.
We consider a total of \NSyst{} energy scale, neutrino flux, neutrino cross-section,
efficiency, initial-state nuclear environment and final-state rescattering systematics 
plus the oscillation parameters (nuisance parameters).
In this Run 5+6 analysis, all the nuisance parameters are marginalized 
~ \cite{CousinsHighland}
.
The nuisance parameters and their effects 
on the extrapolated single $e$-like ring event reconstructed energy spectrum are described in Sec.~\ref{sec:systematics}.
The analysis method is presented in detail in Sec.~\ref{sec:oscfit}
%while the chosen goodness-of-fit test, suitable for the current statistics, is described in Sec.~\ref{sec:gof_test}.
%The performance of the fitter was tested in detail for the Run 5+6 fit using ensembles of toy MC experiments, and
%a summary of the results of these tests is presented in Sec.~\ref{sec:validation}. 

%The 2D confidence regions in $\sin^{2}\overline{\theta}_{23}$ and $|\overline{\Delta m}^{2}_{32}|$ 
%or $|\overline{\Delta m}^{2}_{31}|$ 
%are constructed using the constant-$\Delta \chi^{2}$
%%~ constructed using the Feldman-Cousins 
%method. 
%Systematic errors are included in the oscillation contours using the profiling method.
%%Systematic errors are included in the oscillation contours using the 
%%Cousins-Highland method (Bayesian integration) 
%%~ \cite{CousinsHighland,Kendall}.
%%%Confidence contours constructed with the constant-$\Delta\chi^{2}$ method are also shown for comparison.
%%Both of these methods are summarized in Sec.~\ref{sec:contour_construction}, and
%%the exact methodology used for incorporating systematic errors is explained in detail.

The significance to $\bar{\nu}_e$ appearance was given as a p-value.
Two different test statistics were used: 
the total number of observed events ($N_{obs}$), without using any information from the 
reconstructed energy shape, was used to quote the rate-only p-value,
while $\Delta \chi^2 = \chi(\beta=0) - \chi(\beta=1)$ was used to quote the rate+shape p-value.
More details are given in ~\ref{sec:pvalue}

The results of the measurement
using the Run 5+6 dataset is given in Sec.~\ref{sec:results_bestfit} of this note. 
%while the goodness-of-fit tests are presented in Sec.~\ref{sec:results_gof}.
%The best-fit spectra are shown in Sec.~\ref{sec:results_bestfit} for each of the normal and inverted mass hierarchy assumptions for the 
%Run 5+6 dataset.
%The confidence regions are shown in Sec.~\ref{sec:results_dchisq_contours}. 
%In that section, we also show comparisons of our T2K Run 5+6 results with results from MINOS
%as well as with the expected T2K sensitivity for an integrated exposure of \TotalPot{} POT.

Finally, a summary is given in Sec.~\ref{sec:summary}.  

\clearpage
